The testing infrastructure has been implemented as recomended in RFC 2544, which defines methods for evaluating network performance. 
It consists of a device under test (DUT), connected to a measurement device called \textit{Tester}.\footnote{The hardware used in this testing setup was loaned free of charge for the purposes of this bachelor thesis by Silicon Hill club.}.
In line with the more modern RFC 8219, which states that: ``All tests described SHOULD be performed with bidirectional traffic''~\cite{RFC8219}, the infrastructure is designed to operate with bidirectional traffic. 
This approach ensures more accurate performance measurement under real-world network conditions, as opposed to unidirectional traffic.
The Device Under Test (DUT) and the measurement device are connected using 100Gbit capable cables, preventing any potential bottlenecks in the connection.
The illustration of this hardware setup is shown in fig. \ref{fig:hardware-setup}

\begin{figure}[!htbp]
    \centering
    \includegraphics[width=0.9\linewidth]{images/setup.png}
    \caption{Picture showing hardware setup}
    \label{fig:hardware-setup}
\end{figure}

The Device Under Test (DUT) is the network device being evaluated during testing. 
It is configured with a specific network stack and settings based on measurement scenario 
and serves as the focus of performance and behavior analysis in a controlled test environment. 
The DUT is responsible for processing network traffic and responding to the test conditions set by the measurement device.
Additionally, the electrical power consumption of the DUT is monitored and measured during the tests to assess its energy efficiency under varying loads.
The hardware of DUT is shown in table \ref{tab:hardware_dut}.

\begin{table}[h!]
\centering
\begin{tabular}{|c|c|}
\hline
\textbf{Hardware Component} & \textbf{DUT (Device Under Test)} \\
\hline
\textbf{CPU Model} & 2x Intel(R) Xeon(R) CPU E5-2660 v3 \\
\hline
\textbf{Frequency} & 2.60GHz \\
\hline
\textbf{Cores} & 10 physical cores each \\
\hline
\textbf{Memory (RAM)} & Size, type, speed \\
\hline
\textbf{Network Interface Cards (NIC)} & Mellanox ConnectX-6 Dx (Dual-port) \\
\hline
\end{tabular}
\caption{Hardware details for DUT (Device Under Test)}
\label{tab:hardware_dut}
\end{table}

The Tester (Measurement Device), on the other hand, is responsible in generating the network traffic and capturing the responses from the DUT.
Its physical features are shown in table \ref{tab:hardware_tester}. 

\begin{table}[h!]
\centering
\begin{tabular}{|c|c|}
\hline
\textbf{Hardware Component} & \textbf{Tester (Measurement Device)} \\
\hline
\textbf{CPU Model} & Model, number of cores, frequency \\
\hline
\textbf{Frequency} & \\
\hline
\textbf{Cores} & \\
\hline
\textbf{Memory (RAM)} & Size, type, speed \\
\hline
\textbf{Network Interface Cards (NIC)} & Model, speed, number of ports \\
\hline
\end{tabular}
\caption{Hardware details for Tester (Measurement Device)}
\label{tab:hardware_tester}
\end{table}
